\section{FORMAL RESTATEMENT}
\textbf{Erd\H{o}s \#1; reconstruction of the known counterexample, 2026-09-23.}
A finite set $A$ of positive integers is \emph{dissociated} if its subset
sums are all distinct. The conjecture asks whether an absolute $c>0$
satisfies $\max A\geq c2^{|A|}$ for every such $A$. We prove the negative
answer: for every $\epsilon>0$ there are arbitrarily large $n$ and
dissociated $n$-element sets with $\max A<\epsilon2^n$.
The covering-system question is problem \#2, not this problem.

\section{QUICK LITERATURE AND CONTEXT CHECK}
The live problem page labels this conjecture disproved and reports a Lean
verification. Its exposition by Thomas Bloom, edited 2026-09-03, attributes
the counterexample to GPT-6 Astra. The lattice construction below is a
reconstruction of that already reported result; no novelty is claimed.
I have not independently run the referenced Lean development.
The earlier local attempts concern the correct distinct-sums problem but
do not contain this counterexample.

The source exposition leaves its primitive-lattice approximation step in
a sketch form. We give an explicit integer-matrix argument for precisely
the rational lattices used here. The only standard algebraic input is
Smith normal form for an integer matrix: integral row and column
operations with determinant $\pm1$ reduce a full-column-rank matrix to
a diagonal matrix followed by zero rows. This follows by the Euclidean
algorithm on the entries. No equidistribution theorem is used.

\section{OBJECTIVE AND ATTACK PLAN}
Construct rational lattices in the sum-zero hyperplane that avoid the
open unit cube and have a small normalized covolume. Approximate a slight
dilation by primitive integer lattices using Smith normal form. Their
positive normal vectors give integer weights without short relations;
binary expansion converts these weights into ordinary dissociated sets.
The slight dilation and the dyadic choice of the relation bound are both
essential bookkeeping steps in this argument.

\section{WORK}
\subsection{1. An odd-cycle building block}
Write $V_d=\{x\in\mathbb R^d:\sum_i x_i=0\}$. A pair $(\Lambda,v)$ is
admissible if $\Lambda\subset V_d$ is a lattice of rank $d-1$,
$h:=\sum_i v_i>0$, and
\[
 \Lambda\cap(-1,1)^d=\{0\},\qquad
 (\Lambda+t v)\cap(-1,1)^d\ne\varnothing\ \Longrightarrow\ |t|<1.
 \tag{1}
\]
Take an odd integer $b\geq3$, number coordinates modulo $b$, and set
\[
 (Tz)_i=z_i+\tfrac12z_{i-1},\quad
 \Lambda_1=T(V_b\cap\mathbb Z^b),\quad v_1=Te_0.
\]
Then $h_1=3/2$. The eigenvalue calculation, or expansion of the cyclic
determinant, gives $\det T=1+2^{-b}>0$. In particular, $T$ is invertible.

We check both assertions in (1). Whenever $\|Tz\|_\infty<1$, putting
$M=\|z\|_\infty$ gives $M<1+M/2$, hence $M<2$. If $z$ is integral,
its coordinates belong to $\{-1,0,1\}$. A coordinate equal to $1$ forces
its predecessor to be $-1$, and a coordinate $-1$ forces its predecessor
to be $1$. Any nonzero coordinate would therefore force alternating
signs around an odd cycle, an impossibility. Thus $z=0$, proving the
first assertion even for the larger lattice $T\mathbb Z^b$.

For the second assertion let $z=z'+t e_0$, where $z'\in V_b\cap\mathbb
Z^b$, and suppose $\|Tz\|_\infty<1$. The coordinates $z_1,\ldots,z_{b-1}$
are integral, and $t=\sum_i z_i$. The nonzero entries among these integral
coordinates must form an initial alternating segment: a nonzero entry
forces the preceding entry to have the opposite sign, while a zero
cannot be followed by a nonzero entry. If the entire segment is zero,
$|z_0|<1$ and $t=z_0$. Otherwise change all signs so that $z_1=1$.
The inequality $|1+z_0/2|<1$ implies $z_0<0$. If $z_{b-1}=0$, the
coordinate-zero inequality gives $-1<z_0<0$; the sum of the integral
segment is either $0$ or $1$, so $|t|<1$. If the segment reaches
$b-1$, it has even length, has sum zero, and ends in $-1$. Consequently
$|z_0-1/2|<1$ and $z_0<0$ give $-1/2<t=z_0<0$. This proves (1).

For any admissible pair put $\Gamma=\Lambda+\mathbb Zv$. Taking a basis
of $\Lambda$ and appending $v$ shows
\[
 \operatorname{covol}(\Gamma)
 =\frac{h}{\sqrt d}\operatorname{covol}(\Lambda),\qquad
 \Delta(\Lambda):=\frac{\operatorname{covol}(\Lambda)}{\sqrt d}
 =\frac{\operatorname{covol}(\Gamma)}h.                 \tag{2}
\]
Here $h/\sqrt d$ is the perpendicular component of $v$ relative to
$V_d$. In the building block $\Gamma_1=T\mathbb Z^b$, since
$(V_b\cap\mathbb Z^b)+\mathbb Ze_0=\mathbb Z^b$. Thus
$\operatorname{covol}(\Gamma_1)=1+2^{-b}$.

\subsection{2. Iteration makes the normalized covolume arbitrarily small}
Given an admissible rational pair $(\Lambda_s,v_s)$ in dimension $d$, set
\[
 \Lambda_{s+1}=\left\{(\lambda_j+\beta_jv_s)_{j=0}^{b-1}:
       \lambda_j\in\Lambda_s,\ \beta\in\Lambda_1\right\},
 \qquad v_{s+1}=((v_1)_jv_s)_{j=0}^{b-1}.                 \tag{3}
\]
All these vectors are rational. The sum of block $j$ is $h_s\beta_j$,
so the parametrization is injective, its image lies in $V_{bd}$, and
its rank is $b(d-1)+(b-1)=bd-1$. It is a lattice, as the image of a
product lattice under an injective linear map.

If a vector of $\Lambda_{s+1}$ lies in the open unit cube, the second
part of (1), applied block by block, gives $|\beta_j|<1$. Then
$\beta\in\Lambda_1$ and the first part of (1) imply $\beta=0$.
The remaining $\lambda_j$ vanish by the same first part for $\Lambda_s$.
More generally a vector of $\Lambda_{s+1}+t v_{s+1}$ in the cube has
block coefficients $\gamma=\beta+t v_1$. Their absolute values are
less than one by block admissibility. Admissibility of the building
block then gives $|t|<1$. The new pair is therefore admissible, and
$h_{s+1}=(3/2)h_s$.

To compute volumes, choose a basis $B_s$ of $\Lambda_s$ and use the
block coordinate maps $(x_j,\alpha_j)\mapsto B_sx_j+\alpha_jv_s$.
Each map has determinant of absolute value
$\operatorname{covol}(\Gamma_s)$. The allowed coefficient lattice for
the full-rank lattice $\Gamma_{s+1}$ is
$(\mathbb Z^{d-1})^b\times\Gamma_1$. Hence
\[
 \operatorname{covol}(\Gamma_{s+1})
 =(1+2^{-b})\operatorname{covol}(\Gamma_s)^b.
\]
In dimension $d=b^s$, (2) consequently gives
\[
 \Delta_s=
 \frac{(1+2^{-b})^{1+b+\cdots+b^{s-1}}}{(3/2)^s}.       \tag{4}
\]
Fix $s$ first. As odd $b\to\infty$, the logarithm of the numerator is
at most $(1+b+\cdots+b^{s-1})2^{-b}\to0$. Therefore
$\Delta_s\to(2/3)^s$. Choosing $s$ and then $b$ makes $\Delta_s$ as
small as desired. This order of parameter choices matters.

\subsection{3. Explicit primitive approximation, including dyadic scales}
Let $B$ be a rational $d$ by $m$ basis matrix for one of these lattices,
where $m=d-1$, and let $\Delta=\operatorname{covol}(B\mathbb Z^m)/\sqrt d$.
Choose a positive integer $r$ with $C=rB$ integral. Smith normal form
provides unimodular integer matrices $U,W$ such that
\[
 UCW=D=\begin{pmatrix}\operatorname{diag}(s_1,\ldots,s_m)\\0\end{pmatrix},
 \qquad s_i>0.
\]
Let $E_t$ be the $d$ by $m$ matrix with entries $t s_i$ in position
$(i,i)$, entries $1$ in position $(i+1,i)$, and zero elsewhere.
Define
\[
 A_t=U^{-1}E_tW^{-1}.
\]
These are integer matrices and $A_t/t\to C$. The last $m$ rows of
$E_t$ form a triangular matrix with diagonal entries $1$. Thus its
columns generate a saturated integer lattice: if an integer vector
lies in their real span, its coefficients are integral, as can be
read from these last $m$ coordinates. Unimodular transformations
preserve this property. The signed maximal minors of $A_t$ therefore
give a primitive integral normal vector $u_t$, with
\[
 A_t\mathbb Z^m=u_t^\perp\cap\mathbb Z^d.             \tag{5}
\]
Equivalently, the minor equal to $1$ before the transformations ensures
that the greatest common divisor of all maximal minors is $1$.

Because the columns of $C$ span $V_d$, its signed-minor vector is
$r^m\Delta(1,\ldots,1)$ after choosing its overall sign. Continuity
of the minors gives
\[
 \frac{u_t}{t^m}\longrightarrow r^m\Delta(1,\ldots,1). \tag{6}
\]
In particular all coordinates of $u_t$ are positive for large $t$.

Fix $\delta>0$. For each dyadic $Q=2^\ell$, choose
$t=\lceil(1+\delta)Q/r\rceil$. Then $A_t/Q\to(1+\delta)B$.
For all sufficiently large such $Q$,
\[
 (A_t\mathbb Z^m)\cap(-Q,Q)^d=\{0\}.                 \tag{7}
\]
Indeed, if this failed along an infinite sequence, write a nonzero
short vector as $A_t z_t$. Uniformly bounded left inverses of $A_t/Q$
show that the integer vectors $z_t$ remain bounded. Passing to a
subsequence makes them a fixed nonzero $z$. Taking limits yields
$\|(1+\delta)Bz\|_\infty\leq1$, whereas cube avoidance for $B\mathbb
Z^m$ gives $\|Bz\|_\infty\geq1$, a contradiction. This also explains
why the positive dilation cannot simply be omitted.

It follows from (5)--(7) that the coordinates $u_1,\ldots,u_d$ of
$u_t$ are positive integers satisfying
\[
 \sum_i c_i u_i=0,\quad c_i\in\mathbb Z,\quad |c_i|<Q
 \quad\Longrightarrow\quad c_i=0\text{ for every }i,
 \qquad
 \frac{\max_i u_i}{Q^{d-1}}\longrightarrow(1+\delta)^{d-1}\Delta.
 \tag{8}
\]
This is the required $Q$-dissociation statement at actual powers of two.

\subsection{4. Binary expansion finishes the counterexample}
For the weights in (8), form
\[
 A_Q=\{2^j u_i:1\leq i\leq d,\ 0\leq j<\ell\}.
\]
An equality of two subset sums gives
$\sum_i c_i u_i=0$, where
$c_i=\sum_j\eta_{ij}2^j$, $\eta_{ij}\in\{-1,0,1\}$, and
$|c_i|\leq2^\ell-1=Q-1$. By (8), every $c_i$ vanishes; uniqueness
of binary expansions then says that the two subsets used exactly the
same powers for each $i$. In particular all the displayed elements
are distinct and $A_Q$ is dissociated, with $n=d\ell$. Moreover
\[
 \frac{\max A_Q}{2^n}
 =\frac{(Q/2)\max_i u_i}{Q^d}
 \longrightarrow\frac{(1+\delta)^{d-1}\Delta}{2}.       \tag{9}
\]
Given $\epsilon>0$, use (4) to choose a lattice with $\Delta<\epsilon$.
Then choose $\delta>0$ so small that $(1+\delta)^{d-1}<2$.
The limit in (9) is strictly less than $\epsilon$, and all sufficiently
large $\ell$ give $\max A_Q<\epsilon2^{|A_Q|}$. Their cardinalities
$d\ell$ tend to infinity. This disproves the conjectured uniform
positive constant. \hfill$\square$

\section{ADVERSARIAL VERIFICATION}
The accompanying exact-arithmetic script checks the building block's
integer and single-real-coordinate assertions for $b=3,5,7,9$.
For the latter it intersects the strict rational intervals allowed for
$z_0$, so it checks every possible real value rather than a grid sample.
The proof above handles all odd $b$.

As a concrete check of the primitive approximation when $b=3$, a basis
matrix after clearing denominators is
\[
 C=\begin{pmatrix}1&-1\\1&2\\-2&-1\end{pmatrix}.
\]
The procedure gives
\[
 A_t=\begin{pmatrix}t&-t\\t+1&2t-1\\-2t-1&2-t\end{pmatrix},
 \qquad u_t=(3t^2+t+1,\ 3t^2-t,\ 3t^2).
\]
The script verifies the normal-vector identities and primitivity, then
checks (8) by exhaustive short-relation searches for
$Q=2,4,8,16,32,64,128$ and $t=\lceil3Q/5\rceil$. These finite checks
support the formulas; the limiting argument supplies the theorem for
arbitrary dimension and sufficiently large scales. The script also
demonstrates the failure of the odd-cycle argument for even cycle length.

\section{FINAL STATUS AND REMAINING WORK}
\textbf{SOLVED (reproduction of the known negative answer).}
The argument establishes arbitrarily small $\max A/2^{|A|}$ with
arbitrarily large $|A|$. Its input construction is credited to the
current proof exposition. The matrix approximation, boundary argument,
volume computation, and conversion to subsets have been checked here.
There is no remaining mathematical gap identified in this reconstruction;
the standard Smith normal form theorem is explicitly invoked. The
published formal verification has not been audited here, and this
write-up does not provide an efficient numerical construction at large
dimension or improve the known asymptotic result.

\section{SOURCES AND REPRODUCIBILITY}
T.~F.~Bloom, problem \#1 and its proof exposition,
\url{https://www.erdosproblems.com/1}, accessed 2026-09-23.
The source records the attribution of the construction and the current
problem status. Exact checks for this reconstruction are in
\texttt{verification/solved\_1\_checks.py} and their saved results in
\texttt{verification/solved\_1\_checks.json}.

The estimate concerns completion of this mathematical reconstruction,
not a claim of novelty, a new formal proof, or an improved quantitative
bound.

\section{COMPLETION ESTIMATE}
\textbf{COMPLETION: 100\%}
